\documentclass{beamer}

% Impostazioni lingua e caratteri
\usepackage[italian]{babel}
\usepackage[utf8]{inputenc}
\usepackage[T1]{fontenc}

% Scelta del tema (opzionale, ne ho messo uno pulito)
\usetheme{Madrid}

\title{Disuguaglianze sanitarie in Europa}
\subtitle{Ruolo di reddito, istruzione, welfare e capitale sociale (EHIS)}
\author{Gruppo di lavoro}
\date{\today}

\begin{document}
	
	% SLIDE TITOLO
	\begin{frame}
		\titlepage
	\end{frame}
	
	% SLIDE 1
	\begin{frame}{Contesto e Obiettivo}
		\begin{itemize}
			\item Le disuguaglianze socio-economiche influenzano lo stato di salute
			\item Interesse: capire come i fattori sociali si riflettono sul corpo e sul benessere mentale
			\item Dataset: European Health Interview Survey (EHIS) - Eurostat
			\item Focus su:
			\begin{itemize}
				\item Accesso alle cure (Unmet Needs)
				\item Malattie croniche
				\item Salute mentale
			\end{itemize}
		\end{itemize}
	\end{frame}
	
	% SLIDE 2
	\begin{frame}{Ricerca 1: Welfare e accesso alle cure}
		\textbf{Domanda:}
		\begin{itemize}
			\item In che misura il Welfare State mitiga l'effetto del basso reddito sulla rinuncia alle cure?
		\end{itemize}
		
		\vspace{0.3cm}
		\textbf{Confronto:}
		\begin{itemize}
			\item Modello mediterraneo (Italia, Spagna)
			\item Modello scandinavo (Svezia, Danimarca)
		\end{itemize}
		
		\vspace{0.3cm}
		\textbf{Ipotesi:}
		\begin{itemize}
			\item Nei paesi nordici minori unmet needs a parità di reddito
		\end{itemize}
	\end{frame}
	
	% SLIDE 3
	\begin{frame}{Ricerca 2: Istruzione e salute}
		\textbf{Domanda:}
		\begin{itemize}
			\item In che misura l'istruzione influenza la prevalenza di malattie croniche?
		\end{itemize}
		
		\vspace{0.3cm}
		\textbf{Focus:}
		\begin{itemize}
			\item Diabete, ipertensione
			\item Confronto Nord vs Sud Europa
		\end{itemize}
		
		\vspace{0.3cm}
		\textbf{Ipotesi:}
		\begin{itemize}
			\item L'istruzione agisce come ``scudo protettivo''
			\item Effetto più forte nei sistemi sanitari meno inclusivi
		\end{itemize}
	\end{frame}
	
	% SLIDE 4
	\begin{frame}{Ricerca 3: Salute mentale e capitale sociale}
		\textbf{Domanda:}
		\begin{itemize}
			\item Relazione tra isolamento sociale e depressione (65+)
			\item Il supporto sociale riduce l'impatto della disoccupazione?
		\end{itemize}
		
		\vspace{0.3cm}
		\textbf{Variabili:}
		\begin{itemize}
			\item PHQ-8 (depressione)
			\item Stato occupazionale
			\item Frequenza contatti sociali
		\end{itemize}
		
		\vspace{0.3cm}
		\textbf{Ipotesi:}
		\begin{itemize}
			\item Il supporto sociale ha effetto ``buffer'' (cuscinetto)
		\end{itemize}
	\end{frame}
	
	% SLIDE 5
	\begin{frame}{Ricerca 4: Migranti e disuguaglianze}
		\textbf{Domanda:}
		\begin{itemize}
			\item I migranti (UE e non UE) hanno maggiori bisogni sanitari insoddisfatti rispetto ai nativi?
		\end{itemize}
		
		\vspace{0.3cm}
		\textbf{Approfondimento:}
		\begin{itemize}
			\item Differenza tra:
			\begin{itemize}
				\item Effetto reddito
				\item Barriere non economiche (lingua, accesso, informazione)
			\end{itemize}
		\end{itemize}
		
		\vspace{0.3cm}
		\textbf{Ipotesi:}
		\begin{itemize}
			\item I migranti non UE presentano svantaggi maggiori
		\end{itemize}
	\end{frame}
	
\end{document}